\section{Modern and Contemporary Censures: Reason, Revelation, and Doctrinal Method}

Modern acts make the survey's category discipline especially important. A council can define a dogma; a pope can condemn a specified proposition; the Holy Office can prohibit a book or say that a thesis cannot safely be taught; a doctrinal congregation can notify an author of error or ambiguity; and a canonical act can concern disobedience or schism. These are not interchangeable. Freemasonry, Nazism, political parties, ordinary dissent, and canonical irregularity are therefore not inserted as Christian heresies merely because Catholic authority opposed them. A contextual dossier appears only where the doctrinal object and exact response materially illuminate the boundary.

\begin{heresydossier}{Indifferentism, naturalism, and the religious-liberty distinction}{dossier:indifferentism-naturalism}
\objectfield Nineteenth-century papal documents addressed philosophical naturalism, claims that all religions are equally salvific, state control of the Church, and revolutionary political programs in settings where ``liberty of conscience'' could mean immunity from truth or law as well as civil non-coercion. Gregory XVI's \emph{Mirari vos} of 15 August 1832 and Pius IX's \emph{Quanta cura} with the \emph{Syllabus} of 8 December 1864 are controlling acts. Evidence grade A.

\propositionfield The central religious error is that a person can attain eternal salvation equally well through any profession of religion so long as conduct is morally upright, with conscience treated as autonomous from objective religious and moral truth. Related naturalist propositions make civil society's good independent of religion, subordinate ecclesial authority absolutely to the state, or deny any public duty toward true religion. Those claims must not be conflated with a civil right not to be coerced into religious acts.

\responsefield \emph{Mirari vos} rejects religious indifferentism and an unbounded liberty of conscience; \emph{Quanta cura} condemns specified naturalist and state-absolutist propositions. The attached eighty-item \emph{Syllabus} is an index to earlier acts, organized under headings; each source instrument and context controls its force rather than one invented censure for all eighty. Vatican II, \emph{Dignitatis humanae} 1--2, later teaches civil immunity from coercion in religious matters while expressly retaining the moral duty of persons and societies toward truth.

\aftermathfield Catholic doctrine rejects both coercion into faith and relativism about truth. Consequently, the nineteenth-century acts may not be quoted as though they condemned every modern constitutional liberty, while \emph{Dignitatis humanae} may not be read as teaching that all religions are equally true or salvific.
\end{heresydossier}

\begin{heresydossier}{Hermesianism and methodical doubt}{dossier:hermesianism}
\objectfield Georg Hermes attempted a rational reconstruction of Catholic theology beginning from universal doubt. Gregory XVI's brief \emph{Dum acerbissimas}, 26 September 1835, confirmed prohibition of his works and warned the German hierarchy concerning their system. The works and papal act survive. Evidence grade A/B because the act characterizes a corpus rather than attaching a fixed note to numbered propositions.

\propositionfield The disputed method made positive doubt the beginning of theology, treated reason as the principal norm by which revelation must first be mastered, and generated errors concerning faith and its motives of credibility, God's nature, grace, original sin, and freedom. The issue was not careful rational inquiry but a method in which autonomous demonstration governed assent to revealed mystery.

\responsefield \emph{Dum acerbissimas} confirms the Index judgment and characterizes errors with several terms, including false, rash, tending toward skepticism and indifferentism, erroneous, or savoring of heresy. It does not supply one solemn heresy censure for every Hermesian sentence. Vatican I's \emph{Dei Filius} later taught that faith is a supernatural obedience supported, but not produced, by rational motives of credibility.

\aftermathfield A prolonged conflict over university appointments and clerical formation followed in German-speaking lands. Hermesianism declined as an organized school; Catholic theology's rejection of its method did not abolish apologetics, historical inquiry, or philosophical demonstration.
\end{heresydossier}

\begin{heresydossier}{Bautain, Bonnetty, and the fideist limit}{dossier:fideism-traditionalism}
\objectfield Louis-Eug\`ene Bautain and Augustin Bonnetty reacted against rationalism by assigning human reason too little capacity apart from transmitted revelation or primordial tradition. Bautain signed required propositions beginning on 8 September 1840 and renewed in 1844; Bonnetty accepted propositions imposed by an Index decree of 11 June 1855. Evidence grade A.

\propositionfield The corrected theses denied or endangered reason's ability to establish God's existence and some religious or moral truths, made external signs unable to demonstrate revelation's credibility, or treated tradition as the sole source of intellectual and moral knowledge. The subscriptions instead affirm that reason can know God with certainty, that miracles and prophecy are credible signs, and that reason can lead to faith while mysteries remain known through revelation.

\responsefield The Bautain subscriptions are located at DH 2751--2756 and the Bonnetty propositions at DH 2811--2814. They are corrective professions required of named Catholic thinkers, not conciliar definitions that every use of ``fideism'' identifies a formal heretic. Vatican I's \emph{Dei Filius}, chapter 2 and its canons on revelation, supplied the later dogmatic boundary.

\aftermathfield The episode established a two-sided Catholic answer: rationalism may not absorb faith into autonomous proof, and faith may not deny created reason's real knowledge and the rational credibility of revelation. Appeals to tradition remain legitimate when they do not erase reason or confuse human transmission with the supernatural act of faith.
\end{heresydossier}

\begin{heresydossier}{G\"untherian philosophical theology}{dossier:guntherianism}
\objectfield Anton G\"unther sought to reconstruct Christian mysteries through modern philosophical categories. An Index decree of 8 January 1857 prohibited his works, and Pius IX's letter \emph{Eximiam tuam}, 15 June 1857, explained the doctrinal concern. Evidence grade A/B because the books and prohibition survive but the act expressly did not fix individual propositions and censures.

\propositionfield The concerns involve rational mastery of mysteries; accounts of divine self-consciousness and the Trinity that imperil the unity of divine substance; Christological and anthropological formulations affecting the Incarnation and the soul as substantial form of the body; divine freedom in creation; and a relation of philosophy to faith in which philosophical construction becomes the judge of revelation.

\responsefield \emph{Eximiam tuam} confirms the prohibition but states that the earlier decree had not selected particular propositions or attached determinate theological notes. Its positive boundary protects mystery's intelligibility without claiming that reason can demonstrate or reconstruct the Trinity and Incarnation from its own principles.

\aftermathfield G\"unther and principal followers submitted to the judgment. The act therefore cannot support a personal verdict of obstinate heresy, nor may it be represented as condemning engagement with modern philosophy as such.
\end{heresydossier}

\begin{heresydossier}{Ontologism and immediate knowledge of God}{dossier:ontologism}
\objectfield ``Ontologism'' names a family of claims that the human intellect knows finite realities through an immediate intuition of divine being. The Holy Office examined seven propositions and issued its decree on 18 September 1861. Evidence grade A for the formal list and response.

\propositionfield The propositions make habitual immediate knowledge of God essential to all intellectual knowledge; identify the being first known in things with divine being; deny a real distinction between universal being and God; make other ideas modifications of the idea of God; or describe creatures as present within God in a way that turns creation into divine self-distinction.

\responsefield The Holy Office answered that the seven propositions \emph{tuto tradi non possunt}---they cannot safely be taught (DH 2841--2847). This is a precise lower censure, not an assertion that each thesis was solemnly defined as heretical. Catholic doctrine permits genuine natural knowledge of God from creatures while denying that ordinary human cognition is the beatific or immediate intuition of the divine essence.

\aftermathfield Ontologist language continued to be debated and overlapped with disputes about Rosmini. Similar vocabulary does not prove identical systems; the exact seven propositions and their ``unsafe'' censure control this entry.
\end{heresydossier}

\begin{heresydossier}{Vatican I against pantheism, materialism, and rationalist closure}{dossier:vatican-i-modern-systems}
\objectfield Vatican I's constitution \emph{Dei Filius}, 24 April 1870, answered modern systems by defining positive Catholic doctrine concerning God, creation, revelation, faith, and reason. \emph{Pastor aeternus}, 18 July 1870, separately defined papal primacy and infallibility. Both are dogmatic constitutions of an ecumenical council confirmed by Pius IX. Evidence grade A.

\propositionfield The canons reject denial of the one Creator; materialism; identification of divine and created substance; necessary emanation or divine self-development into finite things; denial of free creation from nothing; denial that God can be known naturally from created things; denial of supernatural revelation, inspiration, miracles, or faith as a supernatural virtue; reduction of assent to scientific demonstration; and a change in dogma's meaning under the name of progress.

\responsefield \emph{Dei Filius}, chapters 1--4 and the corresponding canons, anathematizes the exact contradictory propositions while teaching harmony rather than rivalry between faith and reason. \emph{Pastor aeternus}, chapters 3--4, defines the pope's universal ordinary and immediate jurisdiction and ex cathedra infallibility under stated conditions. The conciliar anathemas are objective doctrinal canons in their historical genre, not automatic findings of present personal culpability.

\aftermathfield The council was suspended during the Franco-Prussian War and never completed its intended ecclesiology. Vatican II later situated primacy within episcopal collegiality without retracting the definitions. ``Rationalism'' and ``pantheism'' remain useful only when tied to the proposition actually meant, not as labels for science, philosophical argument, or every non-Catholic system.
\end{heresydossier}

\begin{heresydossier}{Old Catholic rejection after Vatican I}{dossier:old-catholics}
\objectfield Clergy and laity who rejected the Vatican I definitions organized congresses and eventually a separate episcopal succession. Pius IX addressed the movement in \emph{Etsi multa}, 21 November 1873, and \emph{Graves ac diuturnae}, 23 March 1875. The object includes dogmatic denial, an alternative hierarchy, and schism; these must remain distinct. Evidence grade A.

\propositionfield Movement leaders denied the defined universal primacy and ex cathedra infallibility, alleged that the pope or Catholic Church had fallen into error by defining them, located decisive jurisdiction in community or state-supported structures, and elected Joseph Hubert Reinkens as bishop without papal mandate.

\responsefield \emph{Etsi multa} calls the leaders ``new heretics,'' identifies rejection of infallibility and indefectibility, declares Reinkens's election null, and states canonical penalties. \emph{Graves ac diuturnae} describes the sect as condemned and its adherents as separated and schismatic. The nullity and penalties are juridical acts; the contradiction of \emph{Pastor aeternus} supplies the doctrinal object.

\aftermathfield Old Catholic bodies formed the Union of Utrecht and later developed further doctrinal and liturgical differences. Present descendants are not personally charged with the founders' culpability, and current member churches must be described from their own texts rather than only Pius IX's nineteenth-century polemic.
\end{heresydossier}

\begin{heresydossier}{Rosmini, \emph{Post obitum}, and later clarification}{dossier:rosmini-post-obitum}
\objectfield After Antonio Rosmini's death, the Holy Office decree \emph{Post obitum}, 14 December 1887, condemned and proscribed forty propositions drawn from works published posthumously. The CDF revisited their interpretation in a doctrinal note of 1 July 2001. Evidence grade A.

\propositionfield The extracts concern being and cognition, creation, the soul, the Trinity, Incarnation, Eucharist, original sin, justification, and the beatific state. Read in an idealist or ontologist sense, some could identify created and divine being or dissolve creation and supernatural mystery; Rosmini's broader corpus allowed other readings.

\responsefield \emph{Post obitum} judged the forty propositions not consonant with Catholic truth but did not declare Rosmini personally guilty of formal denial of faith. The 2001 CDF note states that the reasons for doctrinal concern had been superseded because the meanings condemned in 1887 do not correspond to Rosmini's authentic position as clarified by the integral text and later doctrine. It retains the objective force of the decree against the heterodox sense.

\aftermathfield The note removed the earlier obstacles to Rosmini's cause; he was beatified in 2007. This rare authoritative afterlife is a controlling example: a proposition's objectively condemned sense, its extraction from a book, and an author's demonstrable intention are three different judgments.
\end{heresydossier}

\begin{heresydossier}{The tendencies called ``Americanism''}{dossier:americanism}
\objectfield Leo XIII's letter \emph{Testem benevolentiae}, 22 January 1899, addressed Cardinal James Gibbons after controversy over a French presentation of Isaac Hecker. It conditionally rejected tendencies attributed to American Catholics while expressing confidence that the American bishops did not intend them. Evidence grade A.

\propositionfield The rejected tendencies would soften or omit doctrine to attract separated Christians; exalt private liberty over ecclesial authority; make interior inspiration render external guidance superfluous; prefer natural to supernatural virtues and ``active'' to supposedly passive virtues; or depreciate religious vows and institutes as unsuitable to modern life.

\responsefield \emph{Testem benevolentiae} rejects those opinions in the stated senses but does not attach a formal heresy note, establish that Hecker held them, or condemn the American constitutional order. Leo expressly distinguishes the alleged system from American character, laws, and customs that can be praised.

\aftermathfield American bishops denied that the described system represented their Church and submitted to the letter. ``Americanism'' is therefore a cautionary label for specified tendencies and a transatlantic reception dispute, not the name of one proved American heresy.
\end{heresydossier}

\begin{heresydossier}{Modernism, proposition list and constructed system}{dossier:modernism}
\objectfield The Holy Office decree \emph{Lamentabili sane exitu}, 3 July 1907 and approved by Pius X on 4 July, rejected sixty-five propositions drawn from contemporary biblical, historical, and dogmatic writing. Pius X's encyclical \emph{Pascendi dominici gregis}, 8 September 1907, then presented the several roles and claims of a connected Modernist system. Evidence grade A for the acts and principal target writings, while attribution to individual authors remains separate.

\propositionfield The complex restricts reason to phenomena and makes God unknowable; derives faith from immanent religious need or sentiment; severs the Christ of faith from historical knowledge; treats revelation and dogma as expressions of changing consciousness; subordinates Scripture, sacraments, Church, and authority to communal evolution; and permits dogmatic meaning to change substantially. Individual historical-critical conclusions do not constitute that full system.

\responsefield \emph{Lamentabili} reproves and proscribes its sixty-five propositions without publishing a separate theological grade for each. \emph{Pascendi} analyzes and rejects the system, famously calling that constructed whole a ``synthesis of all heresies.'' \emph{Praestantia Scripturae}, 18 November 1907, confirmed the authority and penalties; \emph{Sacrorum antistitum}, 1 September 1910, imposed an anti-Modernist oath as a disciplinary safeguard.

\aftermathfield The oath and older profession were replaced in 1967; disciplinary replacement did not reverse dogmatic teaching. Vatican II's \emph{Dei Verbum} 11--12 requires attention to literary forms, historical setting, and authorial intention while affirming inspiration and scriptural truth. ``Modernism'' must therefore not become a polemical synonym for modern scholarship, doctrinal development, or any disputed contemporary theology.
\end{heresydossier}

\begin{heresydossier}{Atheistic communism as a contextual exclusion}{dossier:atheistic-communism-boundary}
\objectfield Marxist--Leninist movements combined a materialist worldview, revolutionary political organization, and varied state regimes. Pius XI's \emph{Divini Redemptoris}, 19 March 1937, and the Holy Office decree of 1 July 1949 supplied major Catholic responses. This dossier explains a doctrinally condemned ideology that is not ordinarily a Christian heresy under canon 751. Evidence grade A.

\propositionfield The condemned system denies God, spiritual soul, and afterlife; interprets reality through dialectical materialism and class struggle; subordinates the person to collective economic development; and treats religion, morality, family, and property through that reduction. Particular economic reforms or concern for workers do not by themselves constitute this system.

\responsefield \emph{Divini Redemptoris} judges atheistic communism intrinsically incompatible with Christianity while also requiring justice for workers. The 1949 Holy Office responsa prohibited knowing support for the Communist Party and propagation of its materialist anti-Christian doctrine and attached historical canonical consequences to Catholics who professed it. Party discipline, sacramental admission, and doctrinal apostasy are separate objects.

\aftermathfield Many adherents were unbaptized or coerced citizens, so ``heresy'' is ordinarily the wrong canonical category; explicit repudiation of Christianity may instead be apostasy for a baptized person. The fall or transformation of particular regimes did not by itself settle every philosophical proposition, and Catholic social doctrine must not be reduced to a partisan opposite.
\end{heresydossier}

\begin{heresydossier}{Mitigated millenarianism}{dossier:millenarianism-1944}
\objectfield A Holy Office decree considered a mitigated millenarian system in which Christ would reign visibly on earth before the final judgment, whether or not preceded by a resurrection of many just persons. The plenary session answered on 19 July 1944; Pius XII approved the response on 20 July, and the decree was issued on 21 July (AAS 36 [1944] 212; DH 3839). Evidence grade A.

\propositionfield The question was not every theological hope for the fulfillment of history, but a temporal visible reign of Christ before the definitive judgment proposed as an eschatological system. The decree's broadened wording deliberately included a mitigated form.

\responsefield The response says that the system \emph{tuto doceri non posse}---it cannot safely be taught. It does not assign the note ``heresy'' or create a personal penal finding. The positive Catholic confession expects Christ's glorious return, resurrection, judgment, and a kingdom whose consummation is God's gift rather than an intrahistorical political construction.

\aftermathfield The Catechism 675--677 later distinguishes Christian hope from secular or religious pseudo-messianism and cites the rejection of millenarian falsification. The lower censure must remain exact: ``unsafe to teach'' is not interchangeable with a solemn definition of heresy.
\end{heresydossier}

\begin{heresydossier}{Feeneyism and an exclusivist salvation formula}{dossier:feeneyism}
\objectfield A Boston controversy associated with Leonard Feeney interpreted ``outside the Church no salvation'' as requiring explicit visible Catholic membership in every saved person. The Holy Office letter \emph{Suprema haec sacra}, 8 August 1949, was sent to Archbishop Richard Cushing; a separate Holy Office decree of 13 February 1953 excommunicated Feeney. Evidence grade A.

\propositionfield The restrictive thesis excludes salvation wherever explicit canonical incorporation is absent, thereby denying the efficacy of an implicit desire for the Church under grace and confusing the objective necessity of Church and faith with the visible status of every individual. The opposite indifferentist claim---that any religion saves as such---is equally excluded.

\responsefield \emph{Suprema haec sacra} teaches the Church's necessity while distinguishing explicit membership, explicit desire, and implicit desire; any saving desire must be formed within supernatural faith and charity, not reduced to natural goodwill. The 1953 decree states grave disobedience as the disciplinary basis of Feeney's excommunication; it is not a decree declaring him personally a formal heretic.

\aftermathfield Feeney was reconciled before his death without being required to repudiate an invented proposition formula, while communities deriving from the controversy followed different canonical paths. ``Feeneyism'' should name the documented doctrinal error, not convert the disciplinary penalty or every member of a later community into a heresy judgment.
\end{heresydossier}

\begin{heresydossier}{\emph{Humani generis} and dangerous theological reductions}{dossier:humani-generis-errors}
\objectfield Pius XII's encyclical \emph{Humani generis}, 12 August 1950, addressed currents in philosophy, doctrinal development, biblical interpretation, creation, and original sin without naming one organized ``new theology'' or conducting trials of individual theologians. Evidence grade A.

\propositionfield It rejects dogmatic relativism in which opposed concepts become equivalent expressions; restriction of scriptural truth to religious or moral matters; reduction of transubstantiation to symbolism or communal effect; dilution of the Church's identity and necessity; and philosophical systems incompatible with creation and supernatural order. It says that polygenism could not then be embraced because no reconciliation appeared with the revealed and magisterial teaching on original sin.

\responsefield Paragraphs 14--17 protect enduring dogmatic meaning, 22 scriptural truth, 26 Eucharistic conversion, 27 the Church, and 35--37 human origins. Paragraph 36 permits research and discussion concerning evolution of the human body while requiring the immediate creation of spiritual souls. The encyclical grades opinions variously as false, dangerous, or irreconcilable; it does not label every question it discusses heresy.

\aftermathfield The document prompted restrictions and further theological development before Vatican II. It cannot be cited as condemning evolution as such, every ressourcement theologian, or all development in theological language; neither may its doctrinal cautions be treated as silently repealed.
\end{heresydossier}

\begin{heresydossier}{Situation ethics and objective moral order}{dossier:situation-ethics}
\objectfield The Holy Office instruction of 2 February 1956 addressed a ``new morality'' or situation ethics circulating in Catholic education and formation (AAS 48 [1956] 144--145). The object was a method, not every prudent judgment about circumstances. Evidence grade A.

\propositionfield The prohibited system makes the final norm of conduct an individual's concrete, intuitive judgment or personal sincerity rather than objective moral law, so that universal norms can be displaced by the uniqueness of the situation. Circumstances and intention thereby cease merely to specify responsibility and become authority to make an intrinsically disordered choice right.

\responsefield The instruction prohibited teaching or approving the system in universities, seminaries, and religious study houses, and prohibited propagating or defending it in books, lectures, conferences, or any other manner. It named no proponent and did not attach the formal note ``heresy.'' The positive doctrine retains conscience, prudence, circumstances, and subjective culpability within an objective natural and revealed moral order.

\aftermathfield Later Catholic moral theology continued to study conscience and complex cases; \emph{Veritatis splendor} gave a fuller account of intrinsically evil acts and fundamental moral norms. Neither pastoral discernment nor attention to diminished responsibility is identical with the 1956 system.
\end{heresydossier}

\begin{heresydossier}{The postconciliar error inventory and dogmatic meaning}{dossier:postconciliar-error-inventory}
\objectfield The CDF's confidential circular letter to episcopal conferences of 24 July 1966 summarized reported errors and requested episcopal assessment; it was not an ecumenical definition or a sentence against named authors. The CDF declaration \emph{Mysterium Ecclesiae}, dated 24 June 1973, later addressed the Church, infallibility, priesthood, and dogmatic language publicly. Evidence grade A.

\propositionfield The 1966 inventory names restriction of inspiration and inerrancy; dogmatic formulas whose objective meaning changes; the ordinary Magisterium reduced to opinion; relativism; Christ reduced to an ordinary man and miracles or resurrection naturalized; merely symbolic Eucharistic presence without transubstantiation; Mass without sacrifice; penance reduced to social reconciliation; original sin obscured; situation ethics; and ecumenical indifferentism.

\responsefield The circular calls these examples opinions and errors requiring pastoral investigation; its genre does not assign a formal note to each. \emph{Mysterium Ecclesiae} 1--5 defends the Church's concrete continuity, the infallibility of faithful and pastors under their proper conditions, and dogmatic formulas whose expression can develop while their determinate truth remains. Sections 6 and the conclusion defend the distinct ministerial priesthood and theological freedom within the word and Magisterium.

\aftermathfield These acts show that Vatican II did not authorize doctrinal reversal, while also recognizing legitimate research and improved expression. ``Postconciliar error'' is not the name of a single sect, and disagreement over an interpretation must still be tied to a proposition, level of teaching, and competent judgment.
\end{heresydossier}

\begin{heresydossier}{Hans K\"ung and a Catholic-teaching authorization}{dossier:kung-judgment}
\objectfield Hans K\"ung's works prompted prolonged examination and exchanges concerning Church, infallibility, and ministry. The CDF declaration \emph{Christi Ecclesia}, 15 December 1979, records the doctrinal and administrative conclusion. Evidence grade A.

\propositionfield The declaration identifies positions that deny or reduce the Church's infallibility, detach authentic interpretation from the Magisterium, and fail to preserve received teaching concerning the capacity validly to consecrate the Eucharist. It judges the published formulations, not K\"ung's internal forum or every subject in his corpus.

\responsefield \emph{Christi Ecclesia} states that K\"ung had departed from integral Catholic truth in his doctrine and could no longer be considered a Catholic theologian or teach as one. This is a doctrinal and canonical-administrative judgment concerning authorization; the document does not use the procedure or final formula of a personal heresy trial.

\aftermathfield K\"ung remained a priest and continued academic work and dialogue outside a Catholic teaching mandate. The exact consequence---withdrawal of recognition to teach as a Catholic theologian---must not be inflated into excommunication or a judgment on every later reader.
\end{heresydossier}

\begin{heresydossier}{Certain liberation theologies and Marxist reduction}{dossier:liberation-theology}
\objectfield The CDF's \emph{Libertatis nuntius}, 6 August 1984, examined ``certain aspects'' of liberation theology; \emph{Libertatis conscientia}, 22 March 1986, supplied a positive synthesis of Christian freedom and liberation. The plural is essential because the acts do not condemn every theology arising from poverty or oppression. Evidence grade A.

\propositionfield The rejected reduction adopts Marxist analysis uncritically, treats class struggle as a governing hermeneutic and criterion of truth, reduces sin and salvation to political structures, makes orthopraxis displace received faith, legitimates violence, or recasts Christ, Church, Eucharist, and kingdom as instruments of revolutionary struggle.

\responsefield \emph{Libertatis nuntius} warns that accepting the analysis imports its materialist anthropology and total interpretation even where a theologian rejects atheism verbally. It expressly affirms the biblical demand for justice and concern for the poor. \emph{Libertatis conscientia} develops integral liberation from sin together with moral, social, cultural, and political freedom under truth and charity.

\aftermathfield Particular authors later received separate notifications or disciplinary judgments, which must be cited on their own terms. The Church's preferential concern for the poor, opposition to oppression, and Catholic social action are not the censured object, while a good pastoral intention does not make incompatible premises harmless.
\end{heresydossier}

\begin{heresydossier}{The Lefebvre rupture, SSPX status, and sedevacantist claims}{dossier:sspx-sedevacant-boundary}
\objectfield Archbishop Marcel Lefebvre consecrated four bishops without pontifical mandate on 30 June 1988 for the Society of Saint Pius X. John Paul II's \emph{Ecclesia Dei}, 2 July 1988, judged the consecrations a schismatic act. Sedevacantism is a separate family of claims that the Roman See has been vacant for an extended period because recent popes or councils defected from the faith. Evidence grade A for the public acts; no one act defines every later faction.

\propositionfield SSPX disputes have concerned reception of Vatican II, religious liberty, ecumenism, liturgical reform, and the authority of the postconciliar Magisterium; the 1988 act itself concerned unauthorized episcopal consecration and rupture of communion. Sedevacantist arguments add claims about loss of office, indefectibility, or universal deception that differ among groups.

\responsefield \emph{Ecclesia Dei} calls the consecrations schismatic and records the canonical consequences; it does not issue a blanket proposition-level heresy sentence against everyone attached to the older liturgy. Benedict XVI remitted the four bishops' excommunications in 2009, while \emph{Ecclesiae unitatem}, 2 July 2009, stated that the Society lacked canonical status and that doctrinal questions remained unresolved. No universal act has condemned one standardized system called ``sedevacantism.''

\aftermathfield Francis extended faculties for sacramental absolution in \emph{Misericordia et misera} 12 (2016), and the 2017 marriage-faculty letter still described the ``objective persistence'' of the Society's canonical irregularity; those acts did not themselves supply full canonical status or a doctrinal settlement. The DDF warned on 13 May 2026 that announced episcopal ordinations without papal mandate would constitute a schismatic act, and Leo XIV asked the Society to desist on 29 June. The consecrations occurred at \`Ec\^one on 1 July. The signed DDF decree of 2 July states that Alfonso de Galarreta and the four recipients incurred the \emph{latae sententiae} excommunication reserved to the Apostolic See under canon 1387 and also applies canon 1364 \S1 to de Galarreta. It separately states that Bernard Fellay incurred \emph{latae sententiae} excommunication under canon 1364 \S1 through direct participation as co-consecrator and public adherence; it does not describe Fellay's penalty as reserved. The accompanying explanatory note states that SSPX sacred ministers are in schism, are considered schismatics, and are subject to canon 1364 \S1; it does not individually declare every minister excommunicated. The note further states that lay faithful are schismatic and excommunicated when their adherence is formal under the 1996 explanatory note and that the Society's sacramental ministry is illicit, with its Penance and assisted marriages invalid. These controlling status determinations do not turn the SSPX into one proposition-level named heresy or make every attendee a formal adherent. The categories remain distinct: an act may be schismatic, a thesis may contradict primacy or indefectibility, and a person's canonical status may differ from the doctrinal classification of an object.
\end{heresydossier}

\begin{heresydossier}{Priestly ordination reserved to men}{dossier:womens-ordination}
\objectfield Debate over admitting women to ministerial priesthood led John Paul II to issue \emph{Ordinatio sacerdotalis}, 22 May 1994. The CDF answered a doubt on 28 October 1995, and its 1998 commentary on the \emph{Professio fidei} located the teaching within the levels of assent. Evidence grade A.

\propositionfield The contrary thesis says that the Church possesses authority to confer priestly ordination on women, or that the reservation to baptized men is merely a reversible discipline. Questions about women's dignity, gifts, offices not requiring ordination, or prudential structures are not identical with that proposition.

\responsefield \emph{Ordinatio sacerdotalis} declares that the Church has no authority to confer priestly ordination on women and that this judgment is to be held definitively by all the faithful. The 1995 responsum says it belongs to the deposit of faith and has been set forth infallibly by the ordinary and universal Magisterium. The 1998 doctrinal commentary gives it as an example of a second-paragraph truth to be held definitively, not as its example of a newly defined first-paragraph dogma.

\aftermathfield Denial therefore rejects a definitive Catholic doctrine and impairs full communion, but the technical category must remain the one supplied by the 1998 commentary; it should not be called canonical heresy automatically by collapsing all definitive teaching into revealed dogma. Attempts at ordination also carry separate canonical questions and penalties.
\end{heresydossier}

\begin{heresydossier}{Religious pluralism and the uniqueness of Christ}{dossier:religious-pluralism}
\objectfield The CDF declaration \emph{Dominus Iesus}, 6 August 2000, responded to relativist theories of revelation, Christ, salvation, Church, and religions. It reaffirms Catholic doctrine without declaring that every theologian of religions, adherent of another religion, or non-Catholic Christian is a formal heretic. Evidence grade A.

\propositionfield The rejected theses make revelation in Christ limited or incomplete; separate an eternal saving Logos from Jesus; posit a Logos economy beyond the Incarnation or a Spirit economy parallel to Christ; deny Christ's unique and universal mediation; treat religions as complementary salvific ways equal in principle; or imagine Christ's Church as nowhere concretely continuing in history.

\responsefield Sections 5--8 affirm the definitive and complete revelation in Jesus Christ and reserve scriptural inspiration in the proper sense to the canonical books; 9--15 confess the one incarnate Word and unique universal mediation; 16--17 treat the Church's unicity and subsistence; and 20--22 distinguish participation in grace and elements of truth from religious relativism. The declaration describes several propositions as contrary to Catholic faith but does not assign one penal censure to every author.

\aftermathfield The document acknowledges that the Spirit works and elements of truth and goodness are present among peoples and religions, all ordered to Christ's one mediation. Dialogue and respect therefore coexist with proclamation; neither indifferentism nor contempt for persons follows from the dogmatic claim.
\end{heresydossier}

\begin{heresydossier}{Modern notifications: de Mello, Dupuis, and Sobrino}{dossier:modern-notifications}
\objectfield Modern CDF notifications commonly judge writings and require correction without declaring a new named heresy or deciding subjective culpability. Representative cases are the notification on Anthony de Mello of 24 June 1998, the notification on Jacques Dupuis of 24 January 2001, and the notification on Jon Sobrino dated 26 November 2006 and published in 2007. Target texts and notifications survive. Evidence grade A.

\propositionfield The de Mello notification identifies dissolving the personal God, Christ's unique mediation, evil, and final destiny into an undifferentiated spiritual account. The Dupuis notification addresses ambiguous or insufficient formulations concerning Christ's unique mediation, revelation, the Spirit, Church, kingdom, and religions. The Sobrino notification identifies erroneous or dangerous propositions concerning the divinity of Christ, Incarnation, relation of Jesus to the kingdom, self-consciousness, and saving value of his death.

\responsefield De Mello's incompatible positions were judged capable of grave harm. Dupuis was required to assent to clarifying propositions, with the notification expressly not judging subjective intention. Sobrino's propositions were graded erroneous or dangerous, again without a judgment of subjective culpability. Those exact conclusions and document genres govern; none is a universal act creating ``de-Melloism,'' ``Dupuisism,'' or ``Sobrinoism.''

\aftermathfield Notifications can restrict use of a text, require clarification, or identify a doctrinal boundary while leaving other work untouched. They exemplify why ``the CDF investigated an author'' and ``the Church declared that person a heretic'' are historically and canonically different claims.
\end{heresydossier}

\begin{heresydossier}{``Neo-Gnosticism'' and ``neo-Pelagianism'' as analogies}{dossier:neo-gnostic-pelagian-analogy}
\objectfield The CDF letter \emph{Placuit Deo}, issued 22 February 2018, and Francis's exhortation \emph{Gaudete et exsultate}, 19 March 2018, use ancient names to diagnose modern temptations in Christian life. They do not document institutional descent from ancient Gnostics or Pelagius and do not promulgate a new penal catalogue. Evidence grade A.

\propositionfield The ``neo-Gnostic'' analogy concerns a self-enclosed intellectual or subjective salvation that escapes the concrete body, creation, and relationships; the ``neo-Pelagian'' analogy concerns autonomous self-sufficiency, reliance on personal power or observance, and forgetfulness of grace. The analogies select recurrent structural features rather than every proposition of the ancient systems.

\responsefield \emph{Placuit Deo} 3 explicitly limits the comparison to general common features and denies exact identity with the ancient heresies; 6--15 teaches salvation as God's gift in the incarnate Christ, received bodily and communally in the Church. \emph{Gaudete et exsultate} 35--62 applies the terms pastorally to temptations against holiness and grace.

\aftermathfield These labels are useful only with their stated analogical limit. They must not be converted into accusations that a contemporary school or person has been formally adjudged Gnostic or Pelagian, and similarity supplies no historical genealogy.
\end{heresydossier}
